\scanpage{14}

% Left-edge fragments beside the opening paragraph; their textual status is uncertain.
{\reversemarginpar
\margin{\figure{\unsure{\begin{tikzpicture}[x=.035cm,y=-.035cm,line width=.65pt,line cap=round,line join=round]
\draw (0,0)--(1,5)--(12,8) .. controls (23,11) and (24,19) .. (19,20) .. controls (12,22) and (8,12) .. (2,7);
\draw (0,29) .. controls (8,37) and (19,41) .. (24,39) .. controls (33,38) and (31,33) .. (19,30);
\draw[line width=1.1pt] (2,57) .. controls (6,59) and (4,63) .. (2,64);
\end{tikzpicture}}}}}

b) \(\forall\,F\) injectif sur \(\unsure{Y}\), \(\unsure{g|q}^{*}(F)\) est injectif et \(F\unsure{\mathord{-}|\to}g_{*}g^{*}F\) est un épimorphisme (ex: \(g\) inclusion d'un ouvert \(U\hookrightarrow e_{Y}\)).

On prend
\[
\mathbb R\Gamma_{Y}\bigl(\operatorname{Ker}\bigl(\underbrace{C(F}_{\text{rés. inj.}})\to
g_{*}g^{*}\unsure{\mathopen{(}|C}C(F)\bigr)\bigr)\bigr)
\]

Dans le cas général, la difficulté provient du fait que le cône d'un morphisme de complexes (tel que

% Left-edge vertical, horizontal and point fragments beside the middle display.
{\reversemarginpar
\margin{\figure{\unsure{\begin{tikzpicture}[x=.035cm,y=-.035cm,line width=.8pt,line cap=round]
\draw (0,0) .. controls (1,7) and (0,13) .. (0,17);
\draw (0,23)--(13,22);
\fill (0,39) ellipse [x radius=1.2,y radius=2];
\end{tikzpicture}}}}}

\[
F\to g_{*}\deleted{\unsure{C|(}}\bigl(g^{*}(F)\deleted{\unsure{)|C}}\bigr)\deleted{\unsure{g_{*}|*}}
\qquad\text{)}
\]
n'est pas fonctoriel (dans la catégorie dérivée) p. r. à la flèche dont on veut prendre le cône.

% Small clipped left-edge bow beside the Et pourtant paragraph.
{\reversemarginpar
\margin{\figure{\unsure{\begin{tikzpicture}[x=.05cm,y=-.05cm,line width=.8pt,line cap=round]
\draw (0,0) .. controls (7,4) and (9,10) .. (0,13);
\end{tikzpicture}}}}}

Et pourtant, dans le cas particulier actuel il devrait y avoir un choix fonctoriel. ESt\({}_{\unsure{\cdot|\phantom{\cdot}}}\)-ce \unsure{évident|evident} ?
% The preceding low point is at source (658,1206), below-left of the hyphen.
% Detached high stroke at the right of the question, source (1070,1155).
\margin{\unsure{'|\phantom{'}}}

% Small clipped left-edge bow level with the choix fonctoriel line.
{\reversemarginpar
\margin{\figure{\unsure{\begin{tikzpicture}[x=.05cm,y=-.05cm,line width=.8pt,line cap=round]
\draw (0,0) .. controls (7,4) and (9,11) .. (0,15);
\end{tikzpicture}}}}}

Q\uline{uestion pour Illusi}e : Dans sa théorie des déform.tions de schémas en \unsure{groupes|groupe} plats, il tombe sur des \(H^{3}(B_{G}/X,\quad)\) resp. des \(\operatorname{Ext}^{2}(\unsure{K|X};,\quad)\). Peut-on courtcircuiter sa théorie via la théorie (supposée écrite) des Gr-champs - resp. via la théorie des champs de Picard ? J'ai

% Clipped dark edge fragment beside plats, with a separate low point.
{\reversemarginpar
\margin{\figure{\unsure{\begin{tikzpicture}[x=.05cm,y=-.05cm]
\fill (0,0) .. controls (5,2) and (6,6) .. (4,8) -- (1,10) -- (0,7) -- cycle;
\fill (0,22) circle [radius=.6];
\end{tikzpicture}}}}}

% archive sheet number: 13
% A faint straight footer line and scattered surface/edge specks are visible; no textual role is assigned to them.

